\section{Compassion and empathy: what happens if we do nothing?}
\label{sec:compassion}

Every successful health law begins where medicine itself begins, with a person. A
committee may hear statistics from experts all day, but a patient who describes a
long wait for a trial slot, or a parent who describes a child with a rare cancer,
creates an emotional anchor that the numbers alone cannot. Patient stories do more
than move a room: they supply contextual depth that helps a decision maker grasp
the real-world stakes of a bill \cite{novak2020patient}. The honest question such
a story poses is the one a committee cannot avoid. What happens if we do nothing?

Consider the journey the author has documented elsewhere, of a patient moving
through a regulated Physical AI oncology trial from referral to follow-up
\cite{Kawchak2026CancerPatientsJourney}. At each step the patient is asked to
trust a system. Figure 3 traces that path. Compassion asks us to
see the person inside it, and to remember that for many patients the alternative
to a verified, well-run trial is not a safer status quo but no option at all.

\begin{narrativefig}
\node[nfone] (a) at (0,0) {Referral and eligibility};
\node[nftwo] (b) at (3.6,0) {Informed consent};
\node[nfact] (c) at (7.2,0) {Plan verified by gates};
\node[nftwo] (d) at (10.8,0) {Procedure under audit};
\node[nfhope] (e) at (7.2,-2.2) {Transparent outcome};
\draw[nfedge] (a)--(b); \draw[nfedge] (b)--(c); \draw[nfedge] (c)--(d);
\draw[nfedge] (d) to[out=-90,in=0] (e.east);
\end{narrativefig}
\vspace{-0.15cm}
\figcaption{Figure 3. A patient's journey through a verified trial (Mermaid 11 as
colored TikZ): consent and verification precede any action, and the outcome is
recorded transparently.}

The four faces of this appeal are familiar from rare-disease legislation, cancer
screening laws, expanded coverage, and right-to-try initiatives. They are
summarized in \autoref{tab:compassion}, each paired with the cost of inaction that
a patient would name in testimony.

\needspace{9\baselineskip}\begin{table}[H]
\footnotesize
\begin{tabularx}{\textwidth}{@{}L{3.2cm} Y L{4.2cm}@{}}
\toprule
\textbf{Who} & \textbf{What they face} & \textbf{What ``doing nothing'' means}\\
\midrule
Patient with serious illness & Few trials, long waits, narrow eligibility & Time lost that cannot be recovered\\
Family caregiver & Coordinating care across a fragmented system & Burden carried alone, errors unseen\\
Child with a rare cancer & Too few cases for ordinary evidence & A treatable disease left unstudied\\
Person without access & Distance, cost, or coverage barriers & Care that exists but never arrives\\
\bottomrule
\end{tabularx}
\caption{Four patient archetypes and the human cost of inaction.}
\label{tab:compassion}
\end{table}

The point is not to substitute feeling for evidence. It is to make the evidence
memorable by attaching it to the people it protects. H. R. 9510 answers the
compassion appeal not with a slogan but with structure: a verified pathway that a
patient can actually enter, with consent confirmed and every action recorded. The
quiet question follows naturally. If a careful, checked system can be offered, how
did we ever ask patients to accept less?
